% Modelo jerárquico de tres capas en una red de campus
\begin{tikzpicture}[
  equipo/.style={caja, minimum width=2.1cm, minimum height=0.62cm},
  final/.style={caja, fill=white, minimum width=2.1cm, minimum height=0.62cm},
  enlace/.style={FigAcento, thick},
  capa/.style={draw=gray!60, dashed, rounded corners=3pt, fill=gray!6},
  rotulo/.style={font=\footnotesize\bfseries, text=FigAcento, anchor=west}]
  % Bandas de cada capa
  \foreach \y/\nombre in {5.0/Borde, 3.7/Núcleo, 2.4/Distribución, 1.1/Acceso}{
    \draw[capa] (-6.8,\y-0.5) rectangle (6.8,\y+0.5);
    \node[rotulo] at (-6.7,\y) {\nombre};
  }
  % Equipos
  \node[caja, ellipse, minimum width=3cm, minimum height=0.8cm] (isp) at (0,6.35) {Internet (proveedor)};
  \node[equipo, minimum width=3.2cm] (borde) at (0,5.0) {Enrutador y cortafuegos};
  \node[equipo] (n1) at (-1.5,3.7) {Conmutador\\de núcleo 1};
  \node[equipo] (n2) at (1.5,3.7) {Conmutador\\de núcleo 2};
  \node[equipo] (d1) at (-2.7,2.4) {Distribución\\bloque A};
  \node[equipo] (d2) at (2.7,2.4) {Distribución\\bloque B};
  \foreach \i/\x in {1/-3.9, 2/-1.3, 3/1.3, 4/3.9}{
    \node[equipo] (a\i) at (\x,1.1) {Conmutador\\de acceso};
  }
  \node[final] (e1) at (-3.9,-0.3) {Equipos\\cableados};
  \node[final] (e2) at (-1.3,-0.3) {Punto de acceso\\inalámbrico};
  \node[final] (e3) at (1.3,-0.3) {Equipos\\cableados};
  \node[final] (e4) at (3.9,-0.3) {Punto de acceso\\inalámbrico};
  % Enlaces
  \draw[enlace] (isp) -- (borde);
  \draw[enlace] (borde) -- (n1);
  \draw[enlace] (borde) -- (n2);
  \draw[enlace, line width=2pt] (n1) -- (n2);
  \foreach \n in {n1,n2}{
    \foreach \d in {d1,d2}{ \draw[enlace] (\n) -- (\d); }
  }
  \draw[enlace] (d1) -- (a1);
  \draw[enlace] (d1) -- (a2);
  \draw[enlace] (d2) -- (a3);
  \draw[enlace] (d2) -- (a4);
  \foreach \i in {1,...,4}{ \draw[enlace] (a\i) -- (e\i); }
\end{tikzpicture}
